\chapter{Methodology}
\label{ch:methodology}

\section{Research Design}

This project will use a focused literature review to
compare established 5G security mechanisms with proposed
6G capabilities. The analysis will address AI in network
operations, integrated sensing and communication, and
large-scale IoT connectivity.

The review will examine how published security concerns
could affect the tasks of securing, testing, and monitoring
wireless networks. Its findings will consist of a
documented comparison and critical interpretation of
the selected literature.

\section{Source Selection}

The initial source set comprises the five core documents
introduced in the preceding chapters. These include the
selected 5G security specification, the IMT-2030 framework,
the ETSI ISAC report, and two peer-reviewed surveys.

Additional sources will be considered where a comparison
requires further support. Selection will prioritize
peer-reviewed publications and official technical documents
that directly address one of the three capability areas.
The publication or revision date, document version,
and relevance to the research question will be recorded.

Sources will be assessed according to their technical
detail, supporting evidence, and stated limitations.
Broad predictions without a clear supporting argument
will receive limited weight in the analysis.

\section{Data Extraction and Comparison}

Relevant passages will be recorded with their section
or page references. For each capability area, the
comparison matrix will contain:

\begin{enumerate}
    \item The relevant existing 5G security mechanism
          or practitioner task.
    \item The proposed 6G capability and its source.
    \item The security or privacy concern identified
          in the literature.
    \item The potential implication for securing,
          testing, or monitoring the network.
    \item The evidence status, assumptions, and
          remaining uncertainties.
\end{enumerate}

Claims will be classified as specified mechanisms,
framework objectives, proposed solutions, or reported
findings. Implications derived by the project authors
will be explicitly identified as interpretations and
linked to the passages supporting them.

\section{Quality Checks and Limitations}

Both group members will review the completed matrix
and check whether the cited passages support each
entry. Differences in interpretation will be resolved
by returning to the original sources.

The proposed contribution will also be compared with
the selected surveys to identify which aspects they
already address and what additional organization or
analysis our comparison provides.

The focused source selection limits the completeness
of the review. Furthermore, differences in document
purpose and maturity may restrict direct comparisons.
Conclusions about future operational effects will
therefore remain conditional on the assumptions
identified in the literature.